\section{MAP Inference}
\label{sec:map-inference}

Given a Markov Network with score function $f(x, y)$ of the form
\eqref{eq:mn-score}, evaluating the predictor
\eqref{eq:map-prediction} requires solving the \emph{MAP inference}
problem
\begin{equation}
    \hat{y}(x) \;=\; \argmax_{y \in \mathcal{Y}^{N}} f(x, y).
    \label{eq:map-restated}
\end{equation}
Problem~\eqref{eq:map-restated} is NP-hard for general graphs; see,
e.g., \cite{Shimony1994}. Its tractability therefore depends on the
structure of $G$. Two cases are relevant for this thesis: on a linear
chain, exact inference is obtained in polynomial time by dynamic
programming; on the cyclic Sudoku graph, exact inference is
intractable and approximate methods are required.

\subsection{The Viterbi algorithm}
\label{subsec:viterbi}

For the linear-chain MN of Eq.~\eqref{eq:chain-score}, the maximum of
$f(x, y)$ is computed in $O(T \cdot K^2)$ time by the Viterbi
algorithm~\cite{Viterbi1967, Forney1973}, a dynamic-programming
procedure that maintains, for each position $t$ and each state
$j \in \mathcal{Y}$, the score of the best partial labeling that ends
in state $j$ at position $t$:
\begin{align}
    \delta_0(j) &\;=\; f_0(x, j),
    \label{eq:viterbi-init} \\
    \delta_t(j) &\;=\; f_t(x, j)
    \;+\; \max_{i \in \mathcal{Y}}\!\bigl[
        \delta_{t-1}(i) + f_{t-1, t}(i, j)
    \bigr], \quad t = 1, \ldots, T-1.
    \label{eq:viterbi-recursion}
\end{align}
Back-pointers record, at each step, which predecessor $i$ attained
the maximum in Eq.~\eqref{eq:viterbi-recursion}:
\begin{equation}
    b_t(j) \;=\; \argmax_{i \in \mathcal{Y}}\!\bigl[
        \delta_{t-1}(i) + f_{t-1, t}(i, j)
    \bigr].
    \label{eq:viterbi-backpointer}
\end{equation}
The terminal state of the optimal labeling is
$\hat{y}_{T-1} = \argmax_j \delta_{T-1}(j)$, and the remaining
labels are recovered by following the back-pointers:
$\hat{y}_{t-1} = b_t(\hat{y}_t)$ for $t = T-1, \ldots, 1$.
Correctness of Eqs.~\eqref{eq:viterbi-init}--\eqref{eq:viterbi-backpointer}
follows from the optimal-substructure property of the chain: any
prefix of an optimal labeling is optimal among prefixes that end in
the same state. The chain structure ensures that the score
contribution of positions $0, \ldots, t$ depends on positions
$t+1, \ldots, T-1$ only through the choice of $y_t$, so once
$\delta_t(j)$ records the best score ending in state $j$ at
position $t$, the extension to position $t+1$ depends only on
$\delta_t(\cdot)$.

\subsection{The LP relaxation}
\label{subsec:lp-relaxation}

On a graph with cycles --- such as the Sudoku graph --- exact MAP
inference is NP-hard and the dynamic-programming approach of
Section~\ref{subsec:viterbi} no longer applies. A standard
workaround is to \emph{relax} the discrete label assignment to a
continuous one and solve the resulting linear
program~\cite{Werner2007}.

In this thesis the LP relaxation plays a learning role rather than
an inference role: it provides the theoretical foundation for the
LP-relaxed learning loss of Section~\ref{sec:learning-mn} (LP-M3N),
and the per-example dual variables it introduces are jointly
optimized with the model parameters at training time. MAP inference
at evaluation time on cyclic graphs uses a separate approximate
decoder, presented in Section~\ref{subsec:belief-propagation}.
Algorithms that solve the LP exactly --- such as max-sum diffusion
--- exist but are not used here.

Each integer assignment $y_v = a$ is replaced by a continuous variable
$\mu_v(a) \in [0, 1]$ --- a \emph{relaxed indicator} of label $a$ at
node $v$ --- and each edge assignment $(y_v, y_{v'}) = (a, b)$ by
$\mu_{vv'}(a, b) \in [0, 1]$. Natural consistency conditions are
imposed: the indicators at each node sum to one, and the node and
edge indicators marginalize consistently,
\begin{equation}
    \sum_{a \in \mathcal{Y}} \mu_v(a) = 1, \qquad
    \sum_{b \in \mathcal{Y}} \mu_{vv'}(a, b) = \mu_v(a), \qquad
    \sum_{a \in \mathcal{Y}} \mu_{vv'}(a, b) = \mu_{v'}(b).
\end{equation}
Subject to these constraints and $\mu \geq 0$, the LP maximizes
\begin{equation}
    \sum_{v \in V} \sum_{a \in \mathcal{Y}} f_v(x, a)\, \mu_v(a)
    \;+\; \sum_{\{v, v'\} \in E} \sum_{a, b \in \mathcal{Y}}
    f_{vv'}(a, b)\, \mu_{vv'}(a, b).
    \label{eq:lp-relaxation}
\end{equation}
Since every integer assignment is a feasible $\mu$, the LP optimum
upper-bounds the MAP objective.

When the optimal $\mu$ happens to be integer-valued it corresponds
directly to a labeling $y$, which is then guaranteed to be the MAP
solution. This is always the case on chains and trees, where the LP
and Viterbi optima coincide --- a property that underpins the
tightness check used in Chapter~\ref{chap:experiments} to validate
the LP-relaxed \emph{learning} loss of
Section~\ref{sec:learning-mn}. On cyclic graphs the LP may instead
return a \emph{fractional} $\mu$ whose objective strictly exceeds
the MAP optimum, and in such cases a discrete labeling could in
principle be recovered by rounding,
\begin{equation}
    \hat{y}_v \;=\; \argmax_{a \in \mathcal{Y}} \mu_v(a).
    \label{eq:lp-rounding}
\end{equation}
The present thesis does not solve the LP at evaluation time,
however; the relaxation enters only through the LP-relaxed
learning loss derived in Section~\ref{sec:learning-mn}, and the
inference algorithm used at evaluation time is the Belief
Propagation decoder developed next.

\subsection{Loopy Belief Propagation}
\label{subsec:belief-propagation}

The decoder used at evaluation time in this thesis on cyclic graphs
is loopy max-sum Belief Propagation (BP). BP iteratively passes
messages between neighbouring nodes; for each directed edge
$(u, v)$ with $u, v \in V$, the message from $u$ to $v$ at
iteration $\tau$ is
\begin{equation}
    m^{(\tau)}_{u \to v}(b) \;=\;
    \max_{a \in \mathcal{Y}}\!\Biggl[
        f_u(x, a) + f_{uv}(a, b)
        + \!\!\!\sum_{u' \in \mathcal{N}(u) \setminus \{v\}}\!\!\!
        m^{(\tau-1)}_{u' \to u}(a)
    \Biggr],
    \label{eq:bp-update}
\end{equation}
i.e.\ the best score that $u$ can offer $v$ for each label $b$ at
$v$, given the unary at $u$, the pairwise potential, and all
incoming messages at $u$ except the one from $v$. After a fixed
number of iterations, the per-node \emph{belief}
\begin{equation}
    \beta_v(a) \;=\; f_v(x, a)
    \;+\!\!\!\sum_{u \in \mathcal{N}(v)}\!\!\! m_{u \to v}(a),
    \label{eq:bp-belief}
\end{equation}
is read off and the predicted label at node $v$ is
$\hat{y}_v = \argmax_{a \in \mathcal{Y}} \beta_v(a)$.

On tree-structured graphs BP coincides with Viterbi after one
forward pass and is therefore exact. On graphs with cycles ---
``loopy'' BP --- it is no longer guaranteed to converge or to
recover the MAP configuration, but it remains a strong
approximation in practice on the sparse constraint graphs that
arise in Sudoku. BP was chosen over exact LP solvers (such as
max-sum diffusion of Section~\ref{subsec:lp-relaxation}) because
it has a fixed per-example cost and batches naturally over many
test examples on GPU.
